\section{Future Enhancements: Zero-Knowledge Proofs}
	
	Zero-Knowledge Proofs (ZKP) are not required in the current local-only architecture of ViperVault, since no server-side authentication or interaction exists.
	However, once a synchronization backend and OTA update server are introduced, ZKP-based protocols can provide strong security and privacy advantages.
	
	\subsection{Potential Use Cases}
		\begin{itemize}
			\item \textbf{Server Authentication (Sync/Login)}: Instead of sending password-derived hashes to the server, ViperVault can implement a Password-Authenticated Key Exchange (PAKE) with ZK properties, such as OPAQUE.
			This allows the client to prove knowledge of the master password without revealing it or any verifiable hash to the server.
			\item \textbf{Vault Integrity Proofs}: Clients could prove that uploaded vaults are correctly encrypted under the right key, without disclosing the key itself.
		\end{itemize}
	
	\subsection{Advantages}
		\begin{itemize}
			\item No password hashes or secrets ever stored on the server.
			\item Protection against offline brute-force in case of server breach.
			\item Strong alignment with Zero-Knowledge Architecture principles and GDPR/NIS2 minimization of personal data processing.
			\item Clear FOSS trust signal: the server cannot know secrets even in hashed form.
		\end{itemize}
	
	\subsection{Disadvantages}
		\begin{itemize}
			\item Increased implementation complexity, requiring reliable cryptographic libraries and careful protocol design.
			\item Higher performance cost than traditional logins, especially on low-end devices.
		\end{itemize}